\documentclass[a4paper,12pt]{article}
\usepackage{amsmath}
\usepackage[russian]{babel}
\usepackage[utf8]{inputenc}
\usepackage[T2A]{fontenc}
\usepackage{geometry}
\usepackage{listings}
\usepackage{xcolor}
\usepackage{graphicx}
\usepackage{color}
\usepackage{amsopn}
\usepackage{amstext}


\geometry{margin=1.5cm}
% Пример темы в стиле тёмной IDE
\definecolor{intj-bg}{HTML}{2B2B2B}
\definecolor{intj-key}{HTML}{CC7832}
\definecolor{intj-str}{HTML}{6A8759}
\definecolor{intj-comm}{HTML}{808080}
\definecolor{intj-id}{HTML}{A9B7C6}
\definecolor{intj-num}{HTML}{6897BB}

\lstdefinestyle{intellij}{
    language=C++,
    backgroundcolor=\color{intj-bg},
    basicstyle=\ttfamily\small\color{intj-id},
    keywordstyle=\color{intj-key}\bfseries,
    stringstyle=\color{intj-str},
    commentstyle=\itshape\color{intj-comm},
    numberstyle=\color{intj-comm}\tiny,
    stepnumber=1,
    frame=single,
    rulecolor=\color{intj-comm},
    showstringspaces=false,
    breaklines=true,
    tabsize=4
}

\begin{document}

% ---------- Титульный лист ----------
    \begin{titlepage}
        \begin{center}
            \textbf{Федеральное государственное автономное образовательное учреждение высшего образования «Национальный исследовательский университет ИТМО} \\
            \vspace{1cm}

            \textbf{Факультет программной инженерии и компьютерной техники} \\
            \vspace{1cm}

            \Large{\textbf{ОТЧЁТ}} \\
            \Large{\textbf{по задачам: M-P}} \\
            \large{по дисциплине «Алгоритмы и структуры данных»} \\
            \vspace{2cm}
        \end{center}

        \begin{flushright}
            \textbf{Выполнил:} \\
            Ануфриев Андрей Сергеевич \\
            P3219 \\
        \end{flushright}

        \vfill
        \begin{center}
            Санкт-Петербург 2026
        \end{center}
    \end{titlepage}

% ---------- Содержание ----------
    \tableofcontents
    \newpage

% ========== ЗАДАНИЕ 1 ==========


    \section{Задание M. Цивилизация}

    \lstinputlisting[style=intellij, label=lst:taskM]{../tasks/M.cpp}

    \subsection*{Суть решения}
    \textbf{Основная проблема:} найти путь минимальной стоимости между двумя точками на карте с разнородными рельефами, где каждый тип рельефа имеет свою стоимость прохода.\\
    \textbf{Ключевая идея:} это классическая задача поиска кратчайшего пути с взвешенными рёбрами. Поскольку веса рёбер положительны (1 для поля, 2 для леса), применяем алгоритм Дейкстры. Идея состоит в том, чтобы всегда расширять путь из узла с минимальной известной стоимостью, гарантируя, что это оптимальный путь до этого узла.\\
    \begin{enumerate}
        \item \textbf{Выбор оптимального направления}: приоритетная очередь обеспечивает выбор вершины с минимальным расстоянием, избегая неоптимальных путей.
        \item \textbf{Быстрая релаксация}: благодаря max-heap (использованием отрицательных значений) достигаем эффективной работы с $O(log nm)$ на операцию.
        \item \textbf{Восстановление пути}: сохраняем родительскую вершину для каждой достигнутой позиции, что позволяет восстановить весь маршрут в конце.
        \item \textbf{Обработка непроходимых препятствий}: вода блокируется на уровне проверки соседей, неподходящие маршруты просто не рассматриваются.
    \end{enumerate}
    \textbf{Время: O(nm log nm)}\\
    \textbf{Память: O(nm)}


% ========== ЗАДАНИЕ 2 ==========
    \newpage


    \section{Задание N. Свинки-копилки}

    \lstinputlisting[style=intellij, label=lst:taskN]{../tasks/N.cpp}
    \subsection*{Суть решения}
    \textbf{Основная проблема:} найти минимальное количество копилок для разбиения, чтобы получить доступ ко всем деньгам. Это эквивалентно подсчёту циклов в функциональном графе, так как каждый цикл требует разбиения ровно одной копилки.\\
    \textbf{Ключевая идея:} рассматриваем ключи как функциональный граф — для каждой копилки $i$ есть ровно одно исходящее ребро в копилку $keyIn[i]$. Граф из $n$ вершин и $n$ рёбер распадается на несвязные компоненты, каждая из которых содержит ровно один цикл (потому что каждая вершина имеет одно входящее и одно исходящее ребро). Нужно подсчитать количество этих циклов.\\
    \begin{enumerate}
        \item \textbf{Обнаружение цикла}: использование двух состояний (в текущем пути = 1, полностью обработана = 2) позволяет точно определить, когда мы замкнулись в цикл, в отличие от входа в уже обработанную компоненту.
        \item \textbf{Избежание переподсчёта}: каждая компонента рассматривается ровно один раз благодаря маркировке вершин на этап 2 после прохода.
        \item \textbf{Линейная сложность}: каждая вершина посещается максимум дважды, обеспечивая оптимальное время работы даже на больших входах.
    \end{enumerate}
    \textbf{Время: O(n)}\\
    \textbf{Память: O(n)}

% ========== ЗАДАНИЕ 3 ==========
    \newpage


    \section{Задание O. Долой списывание!}

    \lstinputlisting[style=intellij, label=lst:taskO]{../tasks/O.cpp}
    \subsection*{Суть решения}
    \textbf{Основная проблема:} разделить студентов на две группы (списывающие и дающие) так, чтобы каждый обмен записками был между студентами из разных групп.\\
    \textbf{Ключевая идея:} это задача проверки двудольности графа. Граф двудольный тогда и только тогда, когда его можно покрасить в два цвета так, что соседние вершины имеют разные цвета. Это эквивалентно отсутствию циклов нечётной длины. Используем DFS с двухцветной раскраской — если разуем когда-либо присвоить соседу цвет, отличный от требуемого, граф не двудольный.\\
    \begin{enumerate}
        \item \textbf{Правильная раскраска соседей}: формула $3 - c$ гарантирует чередование ровно двух цветов при спуске по DFS.
        \item \textbf{Обнаружение противоречия}: проверка существующего цвета вершины позволяет тут же обнаружить нарушение двудольности без необходимости переопроверки.
        \item \textbf{Связные компоненты}: обработка всех непосещённых вершин гарантирует проверку всех компонент связности графа.
    \end{enumerate}
    \textbf{Время: O(n + m)}\\
    \textbf{Память: O(n + m)}


% ========== ЗАДАНИЕ 4 ==========
    \newpage


    \section{Задание P. Авиаперелёты}

    \lstinputlisting[style=intellij, label=lst:taskP]{../tasks/P.cpp}
    \subsection*{Суть решения}
    \textbf{Основная проблема:} найти минимальный размер топливного бака, достаточный для перелёта из любого города в любой другой (с дозаправками).\\
    \textbf{Ключевая идея:} если зафиксировать размер бака $C$, можно построить граф всех возможных перелётов (рёбра включаются, если расход топлива $\leq C$). Нужна минимальная величина $C$, при которой граф остаётся сильносвязным. Ответ монотонен: если граф сильносвязен при $C$, то сильносвязен и при $C' > C$. Применяем бинарный поиск по $C$ с проверкой сильносвязности.\\
    \begin{enumerate}
        \item \textbf{Проверка сильносвязности}: используем два DFS на прямом и обратном графе. Если из города 0 достижимы все города в прямом графе, и все города могут достичь город 0 в обратном графе, то граф сильносвязен.
        \item \textbf{Монотонность}: поиск оптимального значения гарантирует нахождение минимума благодаря монотонности функции сильносвязности.
        \item \textbf{Эффективность}: бинарный поиск снижает количество проверок с $\text{maxWeight}$ до $\log(\text{maxWeight})$.
    \end{enumerate}
    \textbf{Время: $O(n^2 \log(maxWeight))$}\\
    \textbf{Память: $O(n^2)$}

\end{document}